\documentclass[aspectratio=169,11pt]{beamer}
\usepackage[T1]{fontenc}
\usepackage[utf8]{inputenc}
\usepackage{lmodern}
\usepackage{amsmath,amssymb}
\usepackage{microtype}
\usepackage{hyperref}
\usetheme{default}
\definecolor{ink}{RGB}{26,46,65}
\definecolor{accent}{RGB}{0,111,124}
\definecolor{muted}{RGB}{79,91,103}
\setbeamercolor{normal text}{fg=ink,bg=white}
\setbeamercolor{frametitle}{fg=ink,bg=white}
\setbeamercolor{title}{fg=ink}
\setbeamercolor{structure}{fg=accent}
\setbeamercolor{block title}{fg=white,bg=ink}
\setbeamercolor{block body}{fg=ink,bg=black!4}
\setbeamerfont{frametitle}{size=\Large,series=\bfseries}
\setbeamertemplate{navigation symbols}{}
\setbeamertemplate{footline}{\hfill\color{muted}\scriptsize Aouad, Lykouris \& Zhong (2026) \quad $\cdot$ \quad Repository 1 \quad \insertframenumber\, /\ \inserttotalframenumber\hspace{1.1em}\vspace{0.7em}}
\setbeamersize{text margin left=1.05cm,text margin right=1.05cm}
\newcommand{\highlight}[1]{{\color{accent}\bfseries #1}}
\title{When AI substitutes for human effort}
\subtitle{Aouad, Lykouris \& Zhong (2026) \\ Proposition 2.1 --- a derivative-free reconstruction}
\author{Repository 1 $\mid$ Independent student analysis}
\date{AI and Economic Modeling $\mid$ 2026-II}
\begin{document}

\begin{frame}[plain]
  \vspace{0.8cm}
  {\color{accent}\bfseries REPOSITORY 1 \quad / \quad FIVE-MINUTE RESEARCH BRIEF}\par
  \vspace{1.0cm}
  {\usebeamerfont{title}\color{ink}\huge When AI substitutes\\[0.12cm] for human effort}\par
  \vspace{0.42cm}
  {\large Aouad, Lykouris \& Zhong (2026)}\par
  \vspace{0.22cm}
  {\large Proposition 2.1: a case split, not an FOC}\par
  \vfill
  {\small Source: arXiv:2605.11350v1, Section 2; course Issue \#6}\par
  {\scriptsize github.com/alejandroventuraventurameza-tech/ai-01-aouad}\par
  \vspace{0.5cm}
\end{frame}

\begin{frame}{1. The worker's decision}
  \large
  Skill $s$, effort $e$ and reliable AI assistance $a$ are \highlight{perfect substitutes}:
  \[
    x=s+e+a,\qquad e\geq 0.
  \]
  The worker chooses effort, not total output:
  \begin{block}{Static objective}
    \centering\Large $\displaystyle \max_{e\geq0}\; p(s+e+a)-\gamma e.$
  \end{block}
  \normalsize
  \vspace{0.25cm}
  Assumptions: $p$ is continuous, concave and nondecreasing; $\gamma>0$.
  \vspace{0.3cm}
  \par\highlight{Economic tension:} AI may displace costly effort without reducing contemporaneous production.
\end{frame}

\begin{frame}{2. A critical input level --- with a tie-break}
  \large
  Let $f(x)=p(x)-\gamma x$. Define the \highlight{largest} global maximizer:
  \[
    x^\star=\max\operatorname*{arg\,max}_{x\geq0}\{f(x)\}.
  \]
  \begin{block}{Why is this well defined?}
    \centering
    $\displaystyle \limsup_{x\to\infty}\frac{p(x)}{x}<\gamma
    \quad\Longrightarrow\quad f(x)\to-\infty.$
  \end{block}
  \vspace{0.18cm}
  Continuity then ensures a finite, attained maximum. The \emph{largest}-maximizer convention matters when $f$ has a flat peak.
  \vspace{0.28cm}
  \par\small This is an optimization result: a first-order condition is neither necessary for the proof nor sufficient to handle all corners.
\end{frame}

\begin{frame}{3. Proof by cases (no differentiation)}
  \large
  Put $b=s+a$ and $x=b+e$. Then
  \[
    p(b+e)-\gamma e = f(x)+\gamma b,\qquad x\geq b.
  \]
  \vspace{0.12cm}
  \begin{columns}[T,onlytextwidth]
    \begin{column}{0.47\textwidth}
      \highlight{Case A: $b\leq x^\star$}
      \vspace{0.18cm}
      \par $x^\star$ is feasible and globally optimal. Choose the largest optimizer:
      \[
        e^\star=x^\star-b.
      \]
    \end{column}
    \begin{column}{0.49\textwidth}
      \highlight{Case B: $b>x^\star$}
      \vspace{0.18cm}
      \par Concavity implies $f$ is decreasing beyond its greatest maximizer:
      \[
        \underset{x\geq b}{\arg\max}\, f(x)=\{b\}.
      \]
      Hence $e^\star=0$.
    \end{column}
  \end{columns}
  \vspace{0.18cm}
  \par\normalsize Rigorous inequality: for $y>b>x^\star$, write $b=(1-t)x^\star+ty$, $0<t<1$; concavity and $f(x^\star)\geq f(b)$ yield $f(b)\geq f(y)$.
\end{frame}

\begin{frame}{4. Proposition 2.1 and its meaning}
  \begin{block}{Largest optimal effort and realized gross productivity}
    \centering\Large
    $\displaystyle e^\star(s,a)=(x^\star-s-a)_+$\par\medskip
    $\displaystyle p^\star(s,a)=\max\{p(x^\star),\,p(s+a)\}$
  \end{block}
  \vspace{0.28cm}
  \large
  \begin{itemize}
  \setlength\itemsep{0.23cm}
  \item When $s+a<x^\star$, more AI reduces \emph{effort}, keeping total input at $x^\star$.
  \item Once $s+a\geq x^\star$, effort is zero and gross output is weakly increasing.
  \item This baseline alone \highlight{does not generate a productivity paradox}.
  \end{itemize}
\end{frame}

\begin{frame}{5. What the paper adds --- and one research trap}
  \large
  \textbf{Reading map:} Section 3 (skill dynamics); Section 4 (unreliable AI); Section 5 (AI literacy and polarization). These are \highlight{not reproduced} here.
  \vspace{0.42cm}
  \begin{block}{Do not confuse novelty with an unexamined appendix}
    \normalsize
    Convex effort costs already appear in Appendix D. The later October 2026 v2 also studies forward-looking effort, so myopia is \emph{not} an untouched extension across paper versions.
  \end{block}
  \vspace{0.28cm}
  \par\highlight{Takeaway:} Formal assumptions, tie-breaking and source-version checks matter as much as the algebra.
  \vspace{0.10cm}
  \par{\scriptsize Sources: arXiv:2605.11350v1 (assigned); v2 (October update); course Issue \#6.}
\end{frame}
\end{document}
